% Abstract body only (no sectioning command): PLOS wraps it in \section*,
% PeerJ in the abstract environment.
Gene co-expression analyses identify ``good'' modules by a correlation criterion,
yet standard pipelines detect modules with greedy algorithms that optimize other
quantities and only measure correlation afterwards. We present Genetic Clustering
by Metric (\tool{}), an open-source Python tool that closes this gap by treating
module detection as maximum-likelihood inference under an explicit generative
model and solving it globally. A module is a block of the correlation matrix
following a $q$-factor structure with free loadings; the widely used
absolute-correlation criterion is recovered as the constrained special case in
which $q=1$ and all loadings share a magnitude, so the two sit in one nested
family rather than competing. \tool{} maximizes the likelihood with a memetic
genetic algorithm---a population search hybridized with a greedy refinement that
reassigns genes after the fact, a move the agglomerative clustering at the core of
co-expression pipelines cannot make---and selects the factor rank from the data by
BIC. Against a full set of contemporary module detectors---Leiden, Louvain,
spectral clustering, Markov clustering and WGCNA, each given the true number of
modules or tuned against ground truth---\tool{} attains the best recovery at every
noise level tested (ARI $\FactorAutoHardARI{}$ at $\sigma=1.5$ and
$\FactorAutoExtremeARI{}$ at $\sigma=2.0$, against $\WgcnaOraHardARI{}$ and
$\WgcnaOraExtremeARI{}$ for oracle-tuned WGCNA, \NSeeds{} datasets), and it
isolates unstructured genes far better than any competitor (ARI
$\NoiseGeneMemeticARI{}$ versus $\NoiseGeneLeidenARI{}$ for Leiden). Because the
model is explicit we can also ask when it is wrong: across \MisspecTotal{}
data-generating processes that violate different assumptions the factor model
leads on all \MisspecTotal{}, whereas the constrained special case leads on
\MisspecWins{} and fails outright where modules are not driven by a single
dominant signal---on multi-factor modules it recovers ARI $\MisspecMultiFacGCM{}$
against $\MisspecMultiFacFactor{}$ for the free-loading model. We show this
failure is misspecification rather than search failure, because the optimizer
finds partitions of \emph{higher} likelihood than the ground truth. The
constrained model remains useful as a fast approximation and, unlike the factor
model, its penalty can select the number of modules. On a breast-cancer RNA-seq
dataset the methods prove hard to separate, and we report that the synthetic
advantage does not demonstrably transfer at that sample size. \tool{} depends only on NumPy and SciPy and exposes one swappable-metric
interface with single- and multi-objective modes; the constrained objective
scales to $\ScaleMaxGenes{}$ genes in $\ScaleMaxSeconds{}$~s, while the factor
model is best suited to gene sets of hundreds to low thousands.
